This section is partially formalized. The companion uses Mathlib's actual integer special linear group, its matrices $S,T$, its action on the upper half-plane, and its sets \texttt{fd} and \texttt{fdo}. It does not replace a fundamental domain by a set defined to satisfy the desired conclusions. Steven's correspondence review remains pending.

The \hyperref[modular.generators]{generator proposition} has checked row-operation identities and the complete statement that $S,T$ generate $\operatorname{SL}_2(\mathbb Z)$. The shared row calculations directly match the manuscript. Generation reuses Mathlib's Euclidean reduction for integer matrices of fixed determinant; the library uses integer division, while the prose chooses a nearest integer to reduce the absolute value of the lower entry. The \hyperref[modular.generator-relations]{matrix relations} $S^2=(ST)^3=-I$ and the action $T^n\cdot z=z+n$ are checked. The action of $S$ and the exact effective orders have not yet been exposed in this companion.

The two membership lemmas identify Mathlib's squared-norm inequalities with the manuscript's norm inequalities for the \hyperref[modular.fundamental-region]{closed domain and strict region}. The height lower bound \eqref{eq:modular-height-lower-bound} is checked as $3\leq4y^2$. Orbit reduction proves that every orbit meets the closed domain. Mathlib obtains maximal height through properness of the denominator map, then minimizes the absolute real part among matrices with the selected bottom row; the prose instead makes the finite minimum and the subsequent integer translation explicit.

The companion's final uniqueness lemma requires both representatives to satisfy strict inequalities and concludes equality of points. This is only part of the \hyperref[modular.fundamental-domain]{fundamental-domain theorem}: the prose allows the second representative on the boundary and concludes that the matrix is central. Those stronger statements remain obligations; they have not been weakened in the manuscript.

Further obligations are discreteness and properness; identification of the strict region with the topological interior; nonempty connected interior and its closure; the full denominator bound \eqref{eq:modular-denominator-bound}; centrality and the effective quotient; all boundary equality cases; the exact representative set; and null boundaries and tiling with actual hyperbolic measure. Section 1.7's cusp, stabilizer, and sphere arguments remain separate unimplemented targets.

The machine snapshot is recorded in \texttt{lean/audits/chapter01-release-check-2026-11-12/record.json}. Consult its outcome and source hashes for build and axiom evidence. These checks do not constitute human correspondence approval or complete formalization of this section.
